\section{Related Work}
\label{sec:Related Work}

\textbf{Logical Bug Detection in DBMSs.}
A variety of approaches have been proposed to detect logical bugs in DBMSs~\cite{rigger2020detecting,rigger2020finding,rigger2020testing,song2023testing,tang2023detecting,jiang2024detecting,song2024detecting,hao2023pinolo,deng2025detecting,ba2024keep,jiang2023detecting}.  
For example, \textsc{NoREC}~\cite{rigger2020detecting} transforms an optimizable query into a non-optimizable form and detects semantic inconsistencies by comparing their outputs.  
\textsc{TLP}~\cite{rigger2020finding} divides query predicates into multiple subqueries and verifies that the union of their results is equivalent to the original query.  
\textsc{PQS}~\cite{rigger2020testing} constructs queries expected to retrieve a specific pivot row, while \textsc{DQE}~\cite{song2023testing} detects logical bugs by comparing whether different SQL queries with the same predicate access the same rows in the database.
\textsc{Pinolo}~\cite{hao2023pinolo} generalizes this paradigm through set-level approximation, where predicates are relaxed or restricted to generate over- and under-approximate queries, and correctness is judged by inclusion or containment relations between result sets.
%\textsc{TQS}~\cite{tang2023detecting} decomposes wide tables into smaller ones and uses the base table as ground truth for correctness validation.  
\textsc{EET}~\cite{jiang2024detecting} applies expression-preserving transformations, and \textsc{Radar}~\cite{song2024detecting} compares query results between databases with and without metadata to identify semantic flaws.  
%\textsc{EDC}~\cite{deng2025detecting} detects logical bugs by substituting expressions with precomputed equivalent data and checking for inconsistent results, while \textsc{CODDTest}~\cite{zhang2025constant} leverages constant folding and propagation to generate equivalent queries.  

As discussed in Section~\ref{sec:intro} and Section~\ref{sec:background and motivation}, \toolname\ is orthogonal to these prior efforts.
Most existing techniques design oracles around query structure and relational operators
%(e.g., predicate partitioning, query rewrites, set inclusion/containment, or metadata-aware comparisons)
, and thus excel at exposing bugs in query planning and optimization.
However, they are less effective for logical bugs rooted in built-in SQL functions.
In contrast, \toolname\ targets data-operation-level logical bugs inside built-in SQL function implementations by constructing MRs directly from function properties, enabling function-specific mutations and oracles that remain applicable even when the function is embedded in complex query contexts.


% Test Case Generation
\textbf{DBMS Test-Case Generation.}
A variety of approaches have been proposed for generating diverse test cases for DBMSs, with the aim of improving coverage and revealing potential bugs. 
These techniques~\cite{sqlsmith,zhong2020squirrel,fu2022griffin,jiang2023dynsql,lin2025qtran} typically focus on generating valid and varied SQL queries, but do not specifically target logical bugs. 
For example, \textsc{SQLsmith}~\cite{sqlsmith} is a grammar-based DBMS fuzzer that embeds SQL grammar rules to generate complex SQL queries. 
It uses a random walk approach to explore the SQL syntax and generate a wide range of queries.
\textsc{SQUIRREL}~\cite{zhong2020squirrel} is a mutation-based DBMS fuzzer which introduces an intermediate representation for SQL queries and models dependencies between SQL statements, enabling the generation of queries that contain multiple SQL operations. 
\textsc{Griffin}~\cite{fu2022griffin} uses a grammar-free mutation approach, where it mutates SQL queries based on DBMS state information encapsulated in a metadata graph. 
%\textsc{QTRAN}~\cite{lin2025qtran} is a LLM-based approach that can automatically translate test cases from other DBMSs.
\textsc{DynSQL}~\cite{jiang2023dynsql} takes a dynamic approach by interacting with the DBMS to capture the latest state information, allowing for the incremental generation of valid and complex queries.
These techniques aim to prevent semantic errors and improve the diversity of generated queries. 
In addition to general query generation, some approaches have been specifically designed to aid in bug detection. 
\textsc{SQLRight}~\cite{liang2022detecting} leverages code coverage feedback to enhance test-case generation. 
\textsc{QPG}~\cite{ba2023testing} takes a different approach by recording the query plans covered during DBMS testing and prioritizing the mutation of queries that trigger new query plans. 

These query generation approaches complement \toolname.
They are effective at producing diverse, context-rich queries, while \toolname\ can serve as a lightweight add-on that embeds complex built-in function expressions into these queries and applies function property-driven metamorphic relations as test oracles to detect logical bugs in built-in SQL functions.

